\chapter{结论与展望}\label{chap:conclusion}

\section{主要结论}

本研究建立了含整形板的挤出式3D打印混凝土喷嘴三维CFD模型，针对喷嘴与整形板的4个几何参数，对20个工况进行了系统参数化分析，并构建了五目标Pareto Front分析与NSGA-II优化框架。主要结论如下：

\textbf{（1）模型可靠性}

基于Herschel--Bulkley本构关系建立的非牛顿流体CFD模型通过了网格无关性验证和质量守恒检验（所有工况误差$<1\%$），具备良好的数值可信度。

\textbf{（2）参数重要性排序}

整形板倾斜角对基底接触压力的影响最为显著，整形板长度次之；喷嘴收缩角在显著提升压力的同时还能降低喷嘴壁面剪切应力，具有双重优化效果，而喷嘴竖向高度的影响在统计上并不显著，在实际工程运用中可在工程范围内灵活选取。

\textbf{（3）综合最优配置}

在课题研究构建的几何系统下，C3配置（$TL=60$~mm，$\theta=60°$，$L=20$~mm，$\alpha=15°$）在本研究的离散工况中表现出最优综合性能，TOPSIS综合排名第一。相较于基准工况，其基底压力提升161\%，喷嘴壁面剪切应力降低70\%，压力传递效率$\eta_p=0.560$，混凝土层高$H=6.47$~mm，满足工程层高要求并留有安全余量，是本研究推荐的工程最优方案。

\textbf{（4）能耗与压实的共线性}

在课题研究构建的几何系统下，入口压力与基底压力具有高度共线性。这表明在该类约束流动系统中，提升压实效果伴随泵送能耗的比例增加，两者无法独立优化。这一系统级特性是进行打印参数设计时需要特地权衡的内在约束。

\textbf{（5）NSGA-II优化与边际效益}

经CFD验证的NSGA-II最优解（$TL=70$~mm，$\theta=60°$，$\alpha=14.06°$，$H=5.0$~mm）的基底压力达较C3提升一倍左右，但入口压力同步提升一倍，压力传递效率$\eta_p$仅提升5.9\%。最优解恰好落在层高约束下界，揭示了进一步提升压实效果的边际收益极低，C3是兼顾性能与能耗的工程平衡点。

\section{创新点}

\textbf{（1）建立了含完整整形板结构的三维CFD模型}

本课题首次将顶板、侧挡板与倾斜出料段纳入统一计算域，系统揭示了整形板几何参数对基底接触压力的影响机制。

\textbf{（2）构建了五目标Pareto Front分析与NSGA-II优化框架}

本课题将压力、剪切、效率、均匀性、能耗五个目标进行联合优化，发现入口压力与基底压力存在强共线性，为打印系统的能耗-压实协同设计提供了定量依据。

\section{未来工作展望}

尽管本文建立了基于CFD与代理模型的3D打印混凝土喷嘴优化框架，并获得了若干具有工程意义的结论，但仍存在进一步扩展与完善的空间。

首先，本文采用稳态Herschel--Bulkley模型描述浆体流变行为，尚未考虑材料沉积后的触变性与结构重建过程。未来工作可引入包含结构建立参数的触变性本构模型，并结合瞬态求解方法，对多层打印过程中的早龄期强度演化、层间稳定性以及结构变形行为进行更深入研究。此外，本研究中喷嘴离地高度 $Z_0$ 保持固定，其对接触压力分布与层高形成的独立影响尚未系统量化，后续可围绕不同 $Z_0$ 工况开展进一步参数化分析。

另一方面，本文所有结果均基于数值模拟，尚缺乏实验室打印测试的物理验证。未来可结合实际打印平台，对代表性喷嘴工况开展入口压力、流量及沉积形态等方面的实验对比，以进一步评估CFD模型的预测精度与工程适用性。同时，当前代理模型训练数据规模仍然有限，部分目标函数的预测能力有待提高，后续可通过扩充CFD工况数量与参数覆盖范围，进一步提升NSGA-II优化结果的可靠性。除此之外，未来还可将优化变量从几何参数拓展至材料流变参数，实现喷嘴结构设计与材料配比之间的协同优化。
